\section{Related Work}
\label{sec:related}

\subsection{Single-Image Optical Dehazing and Inpainting}
Early literature in remote sensing cloud removal focused primarily on mono-temporal single-image processing. Methods such as Dark Channel Prior (DCP), multiscale Retinex, and homomorphic filtering attempt to estimate atmospheric transmission maps to remove thin, semi-transparent haze. While effective for cirrus clouds with high transmission coefficients, these mono-temporal methods are fundamentally incapable of recovering surface reflectance beneath opaque, thick cumulus clouds where optical transmission drops to zero ($T(x) \approx 0$). In such cases, classical spatial inpainting algorithms (e.g., Navier-Stokes, Fast Marching, or patch-based exemplar synthesis) hallucinate missing pixels based on adjacent clear-sky textures. However, as cloud patch sizes exceed dozens of meters, spatial inpainting suffers from severe blurring, textural repetition, and complete loss of ground structural veracity \cite{valenta2021cloud}.

\subsection{SAR-to-Optical Image Translation \& Cross-Modal Synthesis}
The advent of deep convolutional neural networks (CNNs) and Generative Adversarial Networks (GANs) \cite{isola2017image} catalyzed research into multi-sensor data fusion. Grohnfeldt et al. \cite{grohnfeldt2018conditional} introduced one of the earliest conditional GAN frameworks (SAR-Opt-cGAN) to translate Sentinel-1 VV/VH polarizations into multispectral Sentinel-2 bands. Schmitt et al. \cite{schmitt2019sen12ms} curated the large-scale SEN12MS dataset, facilitating systematic benchmarking across Earth observation tasks. Meraner et al. \cite{meraner2020cloud} developed DSen2-CR, a deep residual network that fuses same-day Sentinel-1 SAR with cloudy optical imagery. While DSen2-CR demonstrated substantial gains over mono-temporal inpainting, SAR-to-optical translation models frequently suffer from two critical failure modes: (1) \textit{radiometric drift}, where reconstructed vegetation exhibits artificial discoloration due to SAR's inability to directly sense cellular chlorophyll absorption, and (2) \textit{speckle propagation}, where granular SAR noise corrupts high-frequency optical textures.

\subsection{Attention Mechanisms, Cross-Modal Alignment, and Transformers}
To improve cross-modal feature selection, recent architectures have incorporated visual attention mechanisms \cite{vaswani2017attention} and learned spatial registration modules. Xu et al. \cite{xu2022glf} proposed GLF-CR, a global-local fusion network that utilizes self-attention to balance global semantic context with local texture reconstruction. Ebel et al. \cite{ebel2023uncrtaints} introduced UnCRtainTS to formulate multivariate uncertainty estimation in multi-temporal cloud removal. In the cross-attention domain, Hu et al. \cite{hu2025dbcr} introduced DB-CR, employing dual-branch cross-attention to interactively exchange features between SAR and optical streams. Recent studies on multi-sensor resolution discrepancies have explored learned feature alignment frameworks—such as Align-CR \cite{xu2023multimodal} and AFR-CR \cite{zhou2026afrcr}—to bridge spatial coordinate offsets and frequency discrepancies. Recently, Boumahdi et al. \cite{boumahdi2025multimodal} incorporated digital elevation models (DEM) alongside SAR and optical imagery for mountainous cloud removal, validating the importance of topographic relief. However, current paradigms continue to exhibit key limitations:
\begin{enumerate}
    \item \textbf{Simplistic Topographic Ingestion:} Existing DEM-assisted frameworks ingest elevation as an uncalibrated scalar intensity channel, neglecting differential terrain geometry (slope and directional aspect) that physically dictates microwave local incidence angles and cast shadows.
    \item \textbf{Resolution Discrepancy Jitter:} Prior cross-attention architectures predominantly operate on uniform spatial grids (e.g., $10\,\text{m}$ Sentinel-2), suffering from registration jitter and edge blurring when scaling to fine $5.0\,\text{m}$ optical resolutions.
    \item \textbf{Lack of Adaptive Noise Attenuation:} Most high-resolution fusion networks rely on standard deterministic reconstruction objectives, which penalize opaque cloud centers with equal weight as clear ground, causing gradient instability.
\end{enumerate}

\subsection{Diffusion Models vs. Cross-Attention GANs in Remote Sensing}
Recently, Denoising Diffusion Probabilistic Models (DDPMs), such as DiffCR \cite{li2024diffcr}, have emerged for cloud removal. Diffusion models utilize iterative stochastic sampling to synthesize rich high-frequency textures. However, diffusion frameworks require 50 to 100 iterative sampling steps per image patch, resulting in inference latencies exceeding $2.5\,\text{seconds}$ per $256\times 256$ tile. For operational continental-scale or statewide satellite monitoring—where a single full scene contains over $300\,\text{megapixels}$ ($16,000\times 18,000$ pixels)—diffusion inference is computationally prohibitive. In contrast, our proposed \textbf{RelieF-CR} architecture achieves superior structural fidelity (SSIM = 0.963) in a single deterministic forward pass ($\approx 15\,\text{ms}$ per tile on a standard NVIDIA GPU), enabling seamless full-swath reconstruction via memory-mapped streaming.
